\begin{env}[4.3.2]
\label{0.4.3.2-fr}
Sea $(\mathcal{G}_\lambda)$ un sistema inductivo de $\mathcal{B}$-Módulos
y sea $\mathcal{G}=\varinjlim\mathcal{G}_\lambda$ su límite inductivo. Los
homomorfismos canónicos $\mathcal{G}_\lambda\to\mathcal{G}$ definen
$\psi^*(\mathcal{B})$-homomorfismos
$\psi^*(\mathcal{G}_\lambda)\to\psi^*(\mathcal{G})$, que dan un
homomorfismo canónico
\[
  \varinjlim\psi^*(\mathcal{G}_\lambda)
  \longrightarrow
  \psi^*(\mathcal{G}).
\]
Como el tallo en un punto de un límite inductivo de haces es el límite
inductivo de los tallos en ese mismo punto (G, II, 1.11), el homomorfismo
canónico anterior es \emph{biyectivo} (3.7.2). Además, el producto tensorial
conmuta con los límites inductivos de haces, por lo que se tiene un
\emph{isomorfismo canónico funtorial}
\[
  \varinjlim\Psi^*(\mathcal{G}_\lambda)
  \xrightarrow{\sim}
  \Psi^*\bigl(\varinjlim\mathcal{G}_\lambda\bigr)
\]
de $\mathcal{A}$-Módulos.

Por otra parte, para una suma directa finita $\bigoplus_i\mathcal{G}_i$ de
$\mathcal{B}$-Módulos, es claro que
\[
  \psi^*\bigl(\bigoplus_i\mathcal{G}_i\bigr)
  =\bigoplus_i\psi^*(\mathcal{G}_i),
\]
de modo que, al tomar el producto tensorial con $\mathcal{A}_{[\theta]}$,
\[
  \Psi^*\bigl(\bigoplus_i\mathcal{G}_i\bigr)
  =\bigoplus_i\Psi^*(\mathcal{G}_i).
  \tag{4.3.2.1}\label{0.4.3.2.1-fr}
\]
Al pasar al límite inductivo, se deduce, en virtud de lo anterior, que la
igualdad anterior sigue siendo válida para una suma directa
\emph{cualquiera}.
\end{env}
